%-------------------------------------------------------------------------------
%	SECTION TITLE
%-------------------------------------------------------------------------------
\cvsection{Projects \& Contributions}


%-------------------------------------------------------------------------------
%	CONTENT
%-------------------------------------------------------------------------------
\begin{cventries}


%---------------------------------------------------------
  \cventry
    {Wildfire Smoke Detection using Computer Vision} % subtitle
    {Research Project - Early Wildfire Detection} % Project
    {\href{}{}} % Location
    {Nov. 2024} % Location
    {
      \begin{cvitems} % Description(s)
        \item {Developed a method to detect wildfires or initial signs of a fire, capable of running in remote conditions.}
        \item {Preprocessing filters such as Dark Channel Prior, Colour filtering, Corner and Edge detection were utilized to isolate and enhance pixel regions which match the characteristics of fire/smoke.}
        \item {A lightweight YOLOv8 Neural Network was trained on these processed images, enabling inference of fire and smoke within an image. Accuracy of the model saw improvements due to enhancements made by the preprocessing algorithm.}
        \item {Preprocessing pipeline along with Artificial Intelligence model showed improved results on low specification devices such as a Raspberry Pi.}
      \end{cvitems}
    }

%---------------------------------------------------------
  \cventry
    {Neovim Rich Presence Integration Plugin} % subtitle
    {Neopresence} % Project
    {\href{https://github.com/VanillaViking/neopresence}{github.com/VanillaViking/neopresence}} % Location
    {Oct. 2024} % Location
    {
      \begin{cvitems} % Description(s)
        \item {Developed a plugin for Neovim text editor that enables session statistics to be viewed within Discord chat application.}
        \item {Manages Inter process communication using RPC.}
        \item {Efficient asynchronous program, making use of Rust message passing to avoid race conditions.}
        \item {Custom diff calculation algorithm to provide efficient session progress statistics.}
      \end{cvitems}
    }

%---------------------------------------------------------
  \cventry
    {Typing Test Discord Bot} % subtitle
    {Zyenyo} % Project
    {\href{https://zyenyobot.com/}{zyenyobot.com}} % Location
    {Aug. 2022} % Location
    {
      \begin{cvitems} % Description(s)
        \item {Created a bot that allows users to test their typing speed within discord, with functionality such as leaderboards, charts and averages.}
        \item {Used MongoDB aggregation functions to perform statistical analysis on collected typing data.}
        \item {Helped create a type performance rating algorithm, by associating the user's WPM, accuracy \& prompt difficulty}
        \item {Currently active - used by multiple users in a large public server.}
        \item {Utilized Docker to deploy a full stack website along with the bot on a VPS.}
      \end{cvitems}
    }

%---------------------------------------------------------

  \cventry
    {Embedded Systems Project} % Role
    {Proof-Of-Concept Air Conditioning System} % Event
    {\href{https://www.youtube.com/watch?v=NCE8KA36O8I}{youtube link}} % Location
    {May 2023} % Date(s)
    {
      \begin{cvitems} % Description(s)
        \item {Worked in a group to develop a smart air conditioning system that reacts to ambient temperature readings to determine fan operation.}
        \item {Wrote a driver in C for a rotary encoder, which controlled the speed of a motor driving a fan using PWM duty cycle.}
        \item {Utilized the built-in ADC module of the Atmega328P processor to convert analogue values of a thermistor to room temperature values}
        \item {Made extensive use of interrupts, to allow for real-time processing of sensory data.}
      \end{cvitems}
    }


%---------------------------------------------------------
  \cventry
    {Popular Linux Status bar} % Role
    {Polybar Contributions} % Event
    {\href{https://github.com/polybar/polybar}{github.com/polybar}} % Location
    {Aug. 2022} % Date(s)
    {
      \begin{cvitems} % Description(s)
        \item {Fixed a bug that caused the entire status bar to crash when a user provided a certain configuration.}
        \item {Contributed a spacing customisation feature to Polybar (popular status bar for linux), to enable the user to customise token widths in the status bar.}
        \item {Developed an option to prevent the program from setting display manager settings such as "struts", which was reported to cause serious bugs.}
        \item {The pull requests were approved and are currently available to use in the latest version of Polybar.}
      \end{cvitems}
    }

%---------------------------------------------------------
  \cventry
    {On-screen annotation Tool} % Role
    {Gromit-MPX Contributions} % Event
    {\href{https://github.com/bk138/gromit-mpx/pull/180}{github.com/bk138/gromit-mpx}} % Location
    {Sep. 2023} % Date(s)
    {
      \begin{cvitems} % Description(s)
        \item {Added a feature that enables the software to generate strokes from command line input, with customisable characteristics such as colour, thickness and direction/length.}
        \item {Added functionality to save the on-screen notes to a PNG file.}
        \item {The draw-line command was approved and deployed in the 1.5.0 release.}
      \end{cvitems}
    }

%---------------------------------------------------------

  \cventry
    {Speedsolving timer} % Role
    {Virtual Rubik's Cube Model \& Timer} % Event
    {\href{https://cube-timer-theta.vercel.app/}{cube-timer-theta.vercel.app}} % Location
    {Dec. 2023} % Date(s)
    {
      \begin{cvitems} % Description(s)
        \item {Created a speedsolving timer using NextJS that responsively adjusts to desktop or mobile use.}
        \item {Developed a 2D virtual Rubik's cube model, As a requirement of the timer's random scramble tool. This model can be used to double check the accuracy of a user's scramble. }
      \end{cvitems}
    }

%---------------------------------------------------------
\end{cventries}




